\chapter{Podsumowanie i wnioski}

% Zakończenie pracy zwane również Uwagami końcowymi lub Podsumowaniem powinno zawierać ustosunkowanie
% się autora do zadań wskazanych we wstępie do pracy, a w szczególności do celu i zakresu pracy oraz
% porównanie ich z faktycznymi wynikami pracy. Podejście takie umożliwia jasne określenie stopnia
% realizacji założonych celów oraz zwrócenie uwagi na wyniki osiągnięte przez autora w ramach jego
% samodzielnej pracy.

% Integralną częścią pracy są również dodatki, aneksy i załączniki zawierające stworzone w ramach pracy programy, aplikacje i projekty.

\noindent Niniejsza praca poświęcona była analizie sprzętowego generatora liczb 
losowych neoTRNG. W ramach przeprowadzonych badań eksperymentalnych oceniono 
wpływ kluczowych parametrów, metod przetwarzania 
końcowego oraz fizycznego rozmieszczenia elementów w strukturze układu na 
jakość i charakterystykę statystyczną generowanej losowości. 

Na podstawie analizy danych doświadczalnych uzyskanych z pakietów testowych 
NIST SP 800-22, NIST SP 800-90B oraz ENT sformułowano następujące wnioski końcowe:

\section{Wpływ liczby oscylatorów i liczby inwerterów}

Skalowanie zasobów sprzętowych poprzez zwiększanie liczby RO oraz 
liczby inwerterów ma kluczowe znaczenie dla jakości źródła entropii. 

\begin{itemize}
    \item 
    \item 
\end{itemize}

\section{Efektywność przetwarzania końcowego von Neumanna oraz CRC}
Zastosowanie sprzętowych metod post-processingu w module neoTRNG jest warunkiem 
koniecznym do uzyskania wysokiej jakości projektowanego generatora. 

\begin{itemize}
    \item Surowe dane (\textit{FALSE/FALSE}) dyskwalifikują generator w testach 
statystycznych, wykazując silną asymetrię i korelacje lokalne.
    \item
\end{itemize}

\section{Wpływ rozmieszczenia RO}
Analiza rozmieszczenia RO wykazała istotne znaczenie fizycznej lokalizacji 
oscylatorów w strukturze FPGA:
\begin{itemize}
    \item \textbf{Rozmieszczenie zlokalizowane:} 
    \item \textbf{Rozmieszczenie rozproszone:}
\end{itemize}

\section{Zastosowanie zaawansowanych metod przetwarzania końcowego}
Bazowy post-processing neoTRNG może zostać zastąpiony lub uzupełniony przez 
zaawansowane ekstraktory:
\begin{itemize}
    \item \textbf{AES-CTR:}
    \item \textbf{Ekstraktor Toeplitza:} 
\end{itemize}

\section{Podsumowanie końcowe}

Istnieją adekwatne przesłanki do stosowania generatora liczb losowych neoTRNG 
w systemach przetwarzania brzegowego. Świadczą o tym badania przeprowadzone 
i udokumentowane w poprzednich rozdziałach. Wykazano, że kluczem do optymalnego 
zaprojektowania układu TRNG nie jest wyłącznie maksymalizacja liczby 
oscylatorów, lecz balans pomiędzy: dostrojeniem wykorzystania zasobów, rozmieszczeniem 
fizycznym elementów w układzie FPGA, odpowiednim doborem ekstraktora oraz 
rygorystycznym przeprowadzeniem testów.